\chapter{System Testing}
\label{ch:6}


\begin{figure}[H]
\centering
\includegraphics[width=0.8\textwidth]{06_experimental_setups.png}
\caption{Experimental Setup Configuration}
\label{fig:06_experimental_setups}
\end{figure}


\section{6.1 Testing Strategy and Coverage}

\subsection{6.1.1 Test Categories}

\section{6.2 Module 1: Federated Learning Tests}

\subsection{6.2.1 Classification Model Tests}
Test FL-C1: ResNet-18 Forward Pass

\begin{itemize}
  \item **File:** `federated\_learning/tests/test\_model.py`
  \item **Description:** Verify model produces correct output shape and logits
  \item **Input:** Batch of (B, 3, 224, 224) images
  \item **Expected Output:** (B, 3) logits for 3 classes
  \item **Validation:** No NaN/Inf values, output in reasonable range
\end{itemize}

def test\_model\_forward():
    model = BrainTumorResNet(num\_classes=3)
    x = torch.randn(4, 3, 224, 224)
    output = model(x)
    
    assert output.shape == (4, 3), f"Expected (4,3), got \{output.shape\}"
    assert not torch.isnan(output).any(), "Output contains NaN"
    assert not torch.isinf(output).any(), "Output contains Inf"

Test FL-C2: Client Local Training

\begin{itemize}
  \item **Description:** Verify client training loop updates weights and produces decreasing loss
  \item **Dataset:** Toy dataset (100 images, 4 classes)
  \item **Epochs:** 3
  \item **Expected:** Loss decreases, accuracy increases
\end{itemize}

Test FL-C3: Federated Round Execution

\begin{itemize}
  \item **Description:** Execute one complete FL round (broadcast → local training → aggregation)
  \item **Clients:** 3 simulated hospitals
  \item **Expected:** Global model weights updated, test accuracy computed
\end{itemize}


\subsection{6.2.2 Aggregation Strategy Tests}
Test FL-A1: FedAvg Weighted Averaging

\begin{itemize}
  \item **Description:** Verify FedAvg correctly computes weighted average of client weights
  \item **Setup:** 3 clients with dataset sizes [1000, 2000, 3000]
  \item **Validation:** w\_global = (1000/6000)*w1 + (2000/6000)*w2 + (3000/6000)*w3
\end{itemize}

Test FL-A2: FedProx Proximal Regularization

\begin{itemize}
  \item **Description:** Verify proximal term is added to local loss
  \item **Loss Decomposition:** total\_loss = data\_loss + (μ/2)*||w\_local - w\_global||²
  \item **Validation:** Check gradient includes proximal term
\end{itemize}

Test FL-A3: QPSO Particle Updates

\begin{itemize}
  \item **Description:** Verify QPSO correctly updates pbest, gbest, and applies quantum jump
  \item **Validation:**
\end{itemize}

- pbest increases or stays same (never decreases)

- gbest is maximum of all pbests

- New position has bounded magnitude (no explosion)


\section{6.3 Module 2: Segmentation Tests}

\subsection{6.3.1 Data Loading and Preprocessing}
Test SEG-D1: MONAI Transform Pipeline

\begin{itemize}
  \item **File:** `segmentation/test\_preprocessing.py`
  \item **Description:** Verify preprocessing preserves data integrity
  \item **Checks:**
\end{itemize}

- Output shape: (batch, 4, 96, 96, 96) or (4, 128, 128, 128)

- Value range after normalization: typically [-3, +3] (Z-norm)

- No NaN/Inf values

- Orientation is RAS (verified against expected affine matrix)

Test SEG-D2: BraTS Dataset Loading

\begin{itemize}
  \item **File:** `segmentation/test\_data\_loading.py`
  \item **Description:** Load and verify BraTS 2021 dataset structure
  \item **Checks:**
\end{itemize}

- All 4 modalities present (T1, T1ce, T2, FLAIR)

- Ground truth labels present (WT, TC, ET as multi-channel)

- No missing or corrupted files


\subsection{6.3.2 Segmentation Model Tests}
Test SEG-M1: Model Forward Pass

\begin{itemize}
  \item **File:** `segmentation/test\_model.py`
  \item **Input:** (1, 4, 96, 96, 96) - 4-channel 3D volume
  \item **Expected Output:** (1, 3, 96, 96, 96) - 3-channel segmentation
  \item **Validation:** Output in [0, 1] (after sigmoid), no NaN
\end{itemize}

Test SEG-M2: Attention Gate Functionality

\begin{itemize}
  \item **Description:** Verify attention gates suppress background
  \item **Method:** Compare predictions with/without attention
  \item **Expected:** Attention reduces false positives in non-tumor regions
\end{itemize}

Test SEG-I1: Inference with Sliding Window

\begin{itemize}
  \item **File:** `segmentation/test\_inference.py`
  \item **Description:** Test full inference pipeline on sample patient
  \item **Input:** Full 3D volume (\~155×188×155 voxels)
  \item **Expected Output:** Segmentation masks without edge artifacts
  \item **Validation:**
\end{itemize}

- Dice score ≥ 0.75 on held-out test set

- Processing time ≤ 30 seconds per patient

- Tumor Core Dice ≥ 0.80 (clinically important region)


\section{6.4 Module 3: Progression Tests}

\subsection{6.4.1 Mathematical Model Fitting}
Test PROG-M1: Logistic Model Fit

\begin{itemize}
  \item **File:** `progression/test\_growth\_metrics.py`
  \item **Description:** Verify logistic model fits synthetic growth curve correctly
  \item **Setup:** Generate synthetic volume data: V(t) = K/(1 + ((K-V0)/V0)*exp(-r*t))
  \item **Fit:** Recover parameters and verify R² ≥ 0.95
\end{itemize}

Test PROG-M2: Model Selection

\begin{itemize}
  \item **Description:** Test automatic best-fit model selection
  \item **Setup:** 4 synthetic volume trajectories (one each for Logistic, Gompertz, Exponential, Linear)
  \item **Expected:** Best-fit model correctly identifies each trajectory type
\end{itemize}


\subsection{6.4.2 LSTM Hybrid Tests}
Test PROG-L1: LSTM Residual Learning

\begin{itemize}
  \item **File:** `progression/test\_lstm\_training.py`
  \item **Description:** Verify LSTM learns residuals from mathematical baseline
  \item **Setup:**
\end{itemize}

- Train on patient residuals (actual - logistic\_predicted)

- Evaluate on held-out test patients

\begin{itemize}
  \item **Expected:**
\end{itemize}

- Training loss decreases

- Hybrid MAE < Baseline MAE

- Improvement ≥ 5\% for HGG patients

Test PROG-P1: Prediction with Uncertainty

\begin{itemize}
  \item **Description:** Test end-to-end prediction pipeline with confidence intervals
  \item **Input:** Patient with 3-4 historical timepoints
  \item **Output:** 6-month prediction with [lower\_CI, predicted, upper\_CI]
  \item **Validation:**
\end{itemize}

- Confidence interval width is reasonable (5-20\% of prediction)

- Predicted value within 2 standard deviations of held-out test


\section{6.5 Integration and End-to-End Tests}

\subsection{6.5.1 Complete FL Pipeline}
Test INT-FL: 100-Round FL Training

\begin{itemize}
  \item **Scope:** Full federated learning experiment
  \item **Setup:** 3 clients, 100 rounds, 5 local epochs per round
  \item **Metrics Tracked:**
\end{itemize}

- Global test accuracy (should ≥ 95\%)

- Per-client fairness (σ should be minimized)

- Rounds to 80\% accuracy

- Total training time

\begin{itemize}
  \item **Expected Results:**
\end{itemize}

- FedAvg: \~90-98\% final accuracy

- QPSO: \~90-98\% accuracy WITH better fairness (σ ≤ FedAvg σ)


\subsection{6.5.2 Complete Workflow: Classification → Segmentation → Progression}
Test INT-E2E: End-to-End Patient Analysis

\begin{itemize}
  \item **File:** `tests/test\_end\_to\_end.py`
  \item **Workflow:**
\end{itemize}

1. Load sample Glioma patient MRI

2. Run FL Classification (inference only)

3. If Glioma → Run Segmentation

4. Extract tumor volume → Run Progression

5. Return: Diagnosis + Volumetric Analysis + 6-month Prognosis

\begin{itemize}
  \item **Validation:**
\end{itemize}

- All modules execute without errors

- Output types correct (class probabilities, segmentation mask, volume prediction)

- Processing time < 2 minutes total


\section{6.6 Performance and Scalability Tests}

\subsection{6.6.1 Training Performance}

\subsection{6.6.2 Accuracy Benchmarks}
Classification Accuracy:

\begin{itemize}
  \item FedAvg: 98.79\% (Setup 1), 90.56\% (Setup 2, Label Skew)
  \item FedProx: 99.29\% (Setup 1), 93.02\% (Setup 2)
  \item QPSO-FL: 98.43\% (Setup 1), 92.09\% (Setup 2) ← Best fairness
\end{itemize}

Segmentation Metrics (BraTS 2021 Validation Set):

\begin{itemize}
  \item Mean Dice: 0.76 (target ≥ 0.75) ✅
  \item Tumor Core Dice: 0.85 (target ≥ 0.80) ✅
  \item Enhancing Tumor Dice: 0.79 (target ≥ 0.75) ✅
\end{itemize}

Progression Accuracy (MU-Glioma-Post):

\begin{itemize}
  \item HGG MAE (Baseline Logistic): 24,672 mm³
  \item HGG MAE (LSTM Hybrid): 22,728 mm³
  \item **Improvement: 7.88\%** (target ≥ 5\%) ✅
\end{itemize}


\section{6.7 Robustness and Error Handling}

\subsection{6.7.1 Invalid Input Handling}
Test ROB-1: Image Dimension Validation

\begin{itemize}
  \item Input: 256×256 image (wrong size)
  \item Expected: Graceful resize to 224×224 (with warning) OR error with helpful message
\end{itemize}

Test ROB-2: Missing Modalities

\begin{itemize}
  \item Segmentation input: Only T1 provided (should have T1, T1ce, T2, FLAIR)
  \item Expected: Pad missing channels with zeros (with warning)
\end{itemize}

Test ROB-3: Insufficient Timepoints for Progression

\begin{itemize}
  \item Input: Single timepoint (need ≥ 2)
  \item Expected: Error message "Need at least 2 timepoints to fit model"
\end{itemize}


\subsection{6.7.2 Numerical Stability}
Test ROB-4: Large Gradient Values

\begin{itemize}
  \item Input: Very bright MRI slices (causing large gradients)
  \item Expected: Gradient clipping prevents NaN, training continues
\end{itemize}

Test ROB-5: Division by Zero in Logistic Model

\begin{itemize}
  \item Input: V0 = 0 (zero initial volume)
  \item Expected: Handled gracefully (skip patient or use alternative model)
\end{itemize}


\section{6.8 Actual Test Files in Repository}

\subsection{File Inventory}

\section{6.9 Running Tests}

\subsection{Command Line}
\# Run all tests
pytest tests/ -v

\# Run specific module tests
pytest segmentation/test\_inference.py -v
pytest progression/test\_growth\_metrics.py -v

\# Run with coverage report
pytest tests/ --cov=federated\_learning --cov-report=html


\subsection{CI/CD Integration}
Tests automatically run on:

\begin{itemize}
  \item Every git push (GitHub Actions)
  \item Pull requests before merge
  \item Nightly full test suite
\end{itemize}


\section{6.10 Test Summary}
Figure: Experimental Setup Configuration


\section{Chapter Summary}
This chapter documented the testing methodology including unit, integration, and performance testing under IID and non-IID distributions.
